\subsection{Transposition}

In this number, one assumes that X is purely of dimension n.

14.3.1. Let \(\Omega = \det(\mathbf{T}(\mathbf{X})^{*})\) , cf. 7.9.9; this is a vector bundle of rank 1 at each point, with base X, and of class \(\mathbf{C}^{r-1}\) . One has

\[
\Omega = \bigwedge^ {n} \mathrm{T} (\mathrm{X}) ^ {*} = \operatorname{Alt} ^ {n} (\mathrm{T} (\mathrm{X}); \mathrm{K} _ {\mathrm{X}}).
\]

Let M be a vector bundle with base X. One sets

\[
\tilde {\mathbf {M}} = \mathcal {L} (\mathbf {M}; \Omega) = \mathbf {M} ^ {*} \otimes \Omega . ^ {(1)}
\]

If M is of class \(C^{h}\) , with \(h \in N_{K} \cup \{0\}\) and \(h \leqslant r - 1\) , then so is \(\tilde{M}\) . The canonical map from M to \(\tilde{\tilde{M}} = \mathcal{L}(\mathcal{L}(M; \Omega); \Omega)\) is an isomorphism; it is used to identify M with \(\tilde{M}\) .

14.3.2 (“Definition of the transpose”). Let \(k\) be a positive integer \(\leqslant r - 1\) , and let D be a differential operator on X, of type E \(\rightarrow\) F, of order \(\leqslant k\) , and of class C \(^{h}\) , with \(h \in \mathrm{N}_{\mathbf{K}} \cup \{0\}\) and \(k \leqslant h \leqslant r - k\) . There then exists a differential operator \(^{t}\) D on X, of type \(\widetilde{\mathbf{F}} \to \widetilde{\mathbf{E}}\) and of order \(\leqslant k\) , having the following property:

Let \(\xi = (\xi^1, \ldots, \xi^n)\) be a coordinate system in an open set U of X and let \(s = (s_i)_{1 \leqslant i \leqslant e}\) (resp. \(t = (t_j)_{1 \leqslant j \leqslant f}\) ) be a frame of E (resp. of F) on U. Let \((s_i^*)\) (resp. \((t_j^*))\) ) be the frame of E* (resp. of F*) such that \(\langle s_i, s_j^* \rangle = \delta_{ij}\) (resp. \(\langle t_i, t_j^* \rangle = \delta_{ij}\) ), and let \((\tilde{s}_i)\) (resp. \((t_j\) ) be the frame of \(\tilde{E}\) (resp. \(\tilde{F}\) ) on U obtained by forming the tensor product of \((s_i^*)\) (resp. of \((t_j^*)\) ) with the frame \(\omega = d\xi^1 \wedge \cdots \wedge d\xi^n\) of \(\Omega\) . Let \(c_\alpha^{ij}\) ( \(1 \leqslant i \leqslant e\) , \(1 \leqslant j \leqslant f\) , \(|\alpha| \leqslant k\) ) be the functions of class \(C^h\) on U such that one has

\[
\mathrm{D} \left(\sum_ {i} f _ {i} s _ {i}\right) = \sum_ {i, j, \alpha} c _ {\alpha} ^ {i j} \Delta_ {\xi} ^ {\alpha} (f _ {i}). t _ {j}
\]

for any functions \(f_{i}\) in \(\mathcal{C}^{h+k}(\mathbf{U})\) , cf. 14.1.4, (3'). One then has

\[
{ } ^ { t } \mathbf { D } \left( \sum _ { j } g _ { j } \tilde { t } _ { j } \right) = \sum _ { i , j , \alpha } ( - 1 ) ^ { | \alpha | } \Delta _ { \xi } ^ { \alpha } ( c _ { \alpha } ^ { i j } g _ { j } ) \tilde { s } _ { i }
\]

for any functions \(g_{j}\) in \(\mathcal{C}^{h+k}(\mathbf{U})\) .

The preceding property (which must be verified for any \(\xi\) , \(s\) and \(t\) ) determines \(^{t}\) D uniquely. \(^{(1)}\) The operator \(^{t}\) D is called the transpose of D; it is of class \(C^{h-k}\) . If \(h \geqslant 2k\) , the transpose of \(^{t}\) D is defined and equal to D (taking into account the identifications of \(\tilde{E}\) and \(\tilde{F}\) with E and F respectively).

The map \(\mathbf{D} \mapsto {}^t\mathbf{D}\) is K-linear. If \(k = 0\) , i.e., if \(\mathbf{D}\) is a morphism of \(\mathbf{E}\) to \(\mathbf{F}\) , \({}^t\mathbf{D}\) is the morphism \(\mathcal{L}(\mathbf{D}; \mathrm{Id}_{\Omega})\) of \(\tilde{\mathbf{F}}\) to \(\tilde{\mathbf{E}}\) .

The symbol of \({}^{t}D\) is deduced from the symbol of D by means of the isomorphism

\[
\mathsf {T S} ^ {k} (\mathbf {T} (\mathbf {X})) \otimes \mathscr {L} (\mathbf {E}; \mathbf {F}) \rightarrow \mathsf {T S} ^ {k} (\mathbf {T} (\mathbf {X})) \otimes \mathscr {L} (\tilde {\mathbf {F}}; \tilde {\mathbf {E}})
\]

which is the tensor product of the automorphism \((-1)^{k}\) of \(\mathsf{T}\mathsf{S}^{k}(\mathbf{T}(\mathbf{X}))\) and the isomorphism \(u\mapsto\mathcal{L}(u;\mathrm{Id}_{\Omega})\) of \(\mathcal{L}(\mathbf{E};\mathbf{F})\) onto \(\mathcal{L}(\tilde{\mathbf{F}};\tilde{\mathbf{E}})\) .

14.3.3 (“Transpose of a composite”). The notations and hypotheses being those of 14.1.8, suppose that \(D'\) is of class \(C^{h'}\) and \(D''\) of class \(C^{h''}\), with \(h',h'' \in \mathrm{N}_{\mathbf{K}} \cup \{0\}\), \(h' \geqslant k' + 2k''\) and \(h'' \geqslant k'' + 2k'\).

Then the transposes of \,, \,, and \(\circ\) are defined, as is the composite of \(^t\mathbf{D}'\) and \(^t\mathbf{D}''\), and we have

\[
{ } ^ { t } ( \mathbf { D } ^ { \prime \prime } \circ \mathbf { D } ^ { \prime } ) = { } ^ { t } \mathbf { D } ^ { \prime } \circ { } ^ { t } \mathbf { D } ^ { \prime \prime } .
\]

In particular, the hypotheses and notations being those of 14.3.2, and \(f\) being a function of class \(\mathbf{C}^{h'}\) on X with \(2k \leqslant h'\) , one has \({}^t(\mathbf{D} \circ f) = f \circ {}^t\mathbf{D}\) .

Example. --- Take for X an open subset of \(K^{n}\) ; denote by \(\xi^{1}, \ldots, \xi^{n}\) the coordinate functions on X and identify \(\Omega\) with \(K_{x}\) by means of the frame \(\omega = d\xi^{1} \wedge \cdots \wedge d\xi^{n}\) . If \(E = F = K_{x}\) , one has

\[
\tilde {\mathbf {E}} = \tilde {\mathbf {F}} = \mathcal {L} (\mathrm{K}_{\mathbf{X}}; \Omega) = \mathcal {L} (\mathrm{K}_{\mathbf{X}}; \mathrm{K}_{\mathbf{X}}) = \mathrm{K}_{\mathbf{X}},
\]

and the operation \(\mathbf{D} \mapsto {}^{t}\mathbf{D}\) transforms scalar differential operators into scalar differential operators. One has, for example

\[
{ } ^ { t } ( \Delta ^ { \alpha } ) = ( - 1 ) ^ { | \alpha | } \Delta ^ { \alpha } \quad \text { for all } \alpha \in \mathbf {N} ^ { n } .
\]

If \(r\) is infinite, transposition is an anti-automorphism of the algebra \(\mathcal{D}_{\mathbf{x}}^{\infty,r}(\mathrm{K}_{\mathbf{x}},\mathrm{K}_{\mathbf{x}})\) .

14.3.4. Let \(k\) be a positive integer such that K has characteristic 0 or \(>k\) . Let \(\Omega_{1}\) be the vector bundle \(\bigwedge^{n-1}T(X)^{*} = \text{Alt}^{n-1}(T(X); K_{X})\) . The notations and hypotheses being those of 14.3.2, let moreover \(D'\) be a differential operator on X, of type \(\widetilde{F} \to \widetilde{E}\) and of order \(\leqslant k\) . We call any multidifferential operator G (cf. 14.1.11) of type \((E, \widetilde{F}) \to \Omega_{1}\) , of order \(\leqslant k - 1^{(2)}\) and of class \(C^{h}\) (with \(h \geqslant 1\) ), such that one has the identity

\[
\langle \mathrm{D} (u), v \rangle - \langle u, \mathrm{D} ^ {\prime} (v) \rangle = d (\mathrm{G} (u, v))\tag{1}
\]

for every open set U of X and all elements \(u \in \mathcal{S}_{\mathrm{E}}^{k}(\mathrm{U})\) , \(v \in \mathcal{S}_{\mathrm{F}}^{k}(\mathrm{U})\) . Let us specify that, in this formula, \(d\) denotes exterior differentiation (8.3.5) and \(\langle D(u), v \rangle\) (resp.

\(\langle u, \mathrm{D}'(v) \rangle\)) denote the sections of \(\Omega\) over U obtained from the pair \((\mathbf{D}(u), v)\) (resp. the pair \((u, \mathrm{D}'(v)))\)) by the canonical pairing of \(F \times \tilde{F}\) (resp. of \(E \times \tilde{E}\)) into \(\Omega\) . The existence of G implies \(D' = {}^t D\) ; one also says that G is a Green operator for D.

14.3.5. Let D be a differential operator on X, of type E \(\to\) F, of order \(\leqslant k\) , of class \(\mathbf{C}^{h}\) with \(h \in N_{K} \cup \{0\}\) and \(k \leqslant h \leqslant r-k\) . There exists a Green operator of class \(\mathbf{C}^{h-k+1}\) for D in each of the following cases:

\begin{enumerate}
\def\labelenumi{\alph{enumi})}
\item
  \(\mathbf{K} = \mathbf{R}, r = \infty\) and X is paracompact.
\item
  K is of characteristic 0 or \(>k\), X is isomorphic to an open subset of a space \(K^{n}\) , and the bundles E and F are isomorphic to trivial bundles.
\item
  The operator D is of order \(\leqslant1\) (cf. 14.3.7).
\end{enumerate}

14.3.6. We keep the hypotheses and notations of 14.3.3. If G' (resp. G'') is a Green operator for D' (resp. D''), there exists a Green operator H for D'' \(\circ\) D' and only one such that one has

\[
\mathbf {H} (u, v) = \mathbf {G} ^ {\prime \prime} (\mathbf {D} ^ {\prime} (u), v) + \mathbf {G} ^ {\prime} (u, ^ {t} \mathbf {D} ^ {\prime \prime} (v))
\]

for every open set U of X and all elements \(u \in \mathcal{S}_{\mathrm{E}}^{k}(\mathrm{U})\) and \(v \in \mathcal{S}_{\widetilde{\mathrm{G}}}^{k}(\mathrm{U})\) .

14.3.7. We have \(\mathbf{S}^1 (\mathbf{E},\mathbf{F}) = \mathbf{T}(\mathbf{X})\otimes \mathbf{E}^*\otimes \mathbf{F}\) (cf. 14.2.2). On the other hand, the right interior product \((\xi ,\omega)\mapsto i(\xi)\omega\) (cf. A, III, p. 158) defines an isomorphism of \(\mathrm{T(X)}\otimes \Omega\) onto \(\Omega_{1}\) ; by tensor product with the dual \(\Omega^{*}\) of \(\Omega\) , one thus obtains an isomorphism of \(\mathrm{T(X)}\) onto \(\Omega^{*}\otimes \Omega_{1}\) , whence identifications of vector bundles

\[
\begin{array}{r l} \mathrm{S} ^ {1} (\mathrm{E}, \mathrm{F}) & = \Omega^ {*} \otimes \Omega_ {1} \otimes \mathrm{E} ^ {*} \otimes \mathrm{F} = (\mathrm{E} \otimes \mathrm{F} ^ {*} \otimes \Omega) ^ {*} \otimes \Omega_ {1} \\ & = (\mathrm{E} \otimes \tilde {\mathrm{F}}) ^ {*} \otimes \Omega_ {1} = \mathcal {L} (\mathrm{E} \otimes \tilde {\mathrm{F}}; \Omega_ {1}). \end{array}
\]

That said, let D be a differential operator on X, of type E \(\to\) F, of order \(\leqslant1\) , of class \(\mathbf{C}^{h}\) with \(h \in N_{K} \cup \{0\}\) and \(h \leqslant r-1\) . There exists a unique Green operator for D; it is of order 0; it is a morphism of class \(\mathbf{C}^{h}\) from \(\mathbf{E}\otimes\widetilde{\mathbf{F}}\) to \(\Omega_{1}\) , and it is deduced from the symbol of D by the preceding identification.

14.3.8. Suppose that K = R, that X is separated and oriented (10.2.4), and let A be a closed piece (11.1.2) of X. Endow A with the orientation induced by that of X, and \(\partial A\) with the corresponding orientation (11.2.1). Let k be an integer such that \(2k \leqslant r\) , D a differential operator on X, of type E \(\to\) F, of order \(\leqslant k\) and of class \(C^{h}\) . Let G be a Green operator for D. Let U be an open set of X, \(u \in \mathcal{S}_{\mathrm{E}}^{k}(\mathrm{U})\) , \(v \in \mathcal{S}_{\widetilde{\mathrm{F}}}^{k}(\mathrm{U})\) and suppose that the supports of u and of v meet A in a compact set. One then has

\[
\int_ {\mathbf {A}} \langle \mathrm{D} (u), v \rangle - \int_ {\mathbf {A}} \langle u, ^ {t} \mathrm{D} (v) \rangle = \int_ {\partial \mathbf {A}} \mathrm{G} (u, v)
\]

("Green's formula"). In particular, if \(\partial\mathbf{A}=\varnothing\) , one has

\[
\int_ {\mathrm{A}} \langle \mathrm{D} (u), v \rangle = \int_ {\mathrm{A}} \langle u, ^ {t} \mathrm{D} (v) \rangle .
\]


